\section{Permanence while at least \texorpdfstring{$\delta n$}{δn} rows remain}\label{macro}
The goal of this section is Corollary~\ref{macro:cambio}: under the adaptive law $\nu$, except with
probability exponentially small in $n$, every board of the history down to size
$M=\lceil\delta n\rceil$ has a Bethe point that is the unique maximizer and satisfies uniform
diffusion, vacancy, and gap bounds. We prove this first under the auxiliary law $P^{\rm B}$. We
compare the history with a deterministic trajectory (Proposition~\ref{macro:trayectoria}) by means
of observations of rows, columns, fibers, and rectangles (Section~\ref{macro:observaciones}). These
observations allow us to reconstruct the Bethe point (Proposition~\ref{macro:maximo}) and
concentrate around their targets (Proposition~\ref{macro:parada}). Passing from $P^{\rm B}$ to
$\nu$ costs a bounded factor. Comparing a discrete process with the solution of a differential
system, by means of martingales, follows the scheme of Wormald's differential equation
method~\cite{wormald}.

\subsection{Exposure process and comparison trajectory}

Fix $0<\delta<1$. After $s$ choices, write $k=n-s$, $x=k/n$, and $t=s/n$. Under $P^{\rm B}$, we
reveal the uniform random permutation $\sigma$ of the rows one row at a time. Conditional on the
history $\mathcal F_s$ of the choices and the rows already revealed, the next row is uniform among
the $k$ empty rows. If $p$ is the Bethe point of the current board, the conditional probability of
selecting a square $e$ is $p_e/k$. We set $p_e=0$ on squares outside the support. The discrete
indices run from $1$ to $n$. The physical centers $r_i=(i-\tfrac12)/n$ and $c_j=(j-\tfrac12)/n$
place the diagonal $D(e)=i-j$ at $r_i-c_j$ and the antidiagonal $A(e)=i+j-1$ at $r_i+c_j$.
Let $X_f$ denote the availability indicator of a line and $Y_e=\prod_{f\ni e}X_f$ that of a square.
As in Section~\ref{leyes}, $\ell_f$ and $h_f$ are the mass and vacancy of a diagonal or antidiagonal
$f$ at the point $p$.

\begin{lema}[Loss of a square]\label{macro:casilla}
The conditional probability of occupying a given available row or column is $1/k$. The probability of
occupying a given available diagonal $f$ is $\ell_f/k$. For $e=(i,j)$ we have
\[
 \E_{\rm B}[Y_e(s+1)-Y_e(s)\mid\mathcal F_s]
 =-\frac{Y_e(s)}{k}\bigl(2+\ell_{D(e)}+\ell_{A(e)}-3p_e\bigr).
\]
\end{lema}
\begin{proof}
Summing $p_e/k$ over the corresponding line gives the first two formulas. If $Y_e=1$, the four
lines through $e$ intersect only at that square. The sum of their masses counts $e$ four times, and
it must be counted once. If $Y_e=0$, no choice can make the square available again.
\end{proof}

On the unit square we use the profile
\[
 \nu_1(r,c)=b_0r_R(r)r_C(c)u_D(r-c)u_A(r+c).
\]
Its row and column marginals are both one. For $d\in[-1,1]$ and $a'\in[0,2]$, its projections onto
the two diagonal families, that is, the masses of its fibers, are
\[
 m_D(d)=\int_{\{0\le r,r-d\le1\}}\nu_1(r,r-d)\,dr,\qquad
 m_A(a')=\int_{\{0\le r,a'-r\le1\}}\nu_1(r,a'-r)\,dr .
\]
Here and below, the notation under the integral sign requires both coordinates to belong to $[0,1]$.
The vacancies are $u_f=1-m_f$. We use the bounds $c_\nu\le\nu_1\le C_\nu$ and $u_f\ge\kappa_\nu>0$
and the Lipschitz regularity of the full-board profile \cite[Theorems~1.1--1.3, Proposition~5.2 and
Theorem~7.7]{perfil}. The fiber masses are Lipschitz in their parameter. Indeed, shifting the
parameter by $\epsilon$ moves each endpoint of integration by at most $|\epsilon|$ and changes the
integrand on the overlap by at most $L_\nu|\epsilon|$.

The following proposition gives the deterministic targets against which we compare the process when
a fraction $x$ of the rows remain. For a diagonal $f$, $U_f(x)$ is the target for its availability
indicator and $\ell_f^*(x)$ is the target for its mass. At a point $(r,c)$, $J_x$ is the target for
the availability indicator of the square and $v_x/n$ is the target for the Bethe weight of an
available square.

\begin{proposicion}[Comparison trajectory]\label{macro:trayectoria}
For a diagonal line $f$, a point $(r,c)$, and $0\le x\le1$, write
\begin{gather*}
 U_f(x)=u_f+xm_f,\qquad \ell_f^*(x)=\frac{xm_f}{U_f(x)},\\
 J_x=x^2U_D(x)U_A(x),\qquad
 W_x=\frac{\nu_1}{U_D(x)U_A(x)}.
\end{gather*}
Subscripts $D,A$ without arguments indicate evaluation at $r-c,r+c$, respectively. Since $u_f\ge
\kappa_\nu$, $W_x$ is continuous down to $W_0=\nu_1/(u_Du_A)=b_0r_Rr_C$, the weight used in
Section~\ref{meso:regularidad}. We have $J_1=1$. For $0<x\le1$, set $v_x=W_x/x$. Then
$J_xv_x=x\nu_1$, both marginals of $J_xv_x$ are equal to $x$, and
\begin{equation}\label{macro:derivadas}
 \partial_t U_f=-\frac{\ell_f^*}{x}U_f,\qquad
 \partial_t\log J_x=-\frac{2+\ell_D^*+\ell_A^*}{x},\qquad
 1-\ell_f^*=\frac{u_f}{U_f}\ge\kappa_\nu .
\end{equation}
In particular,
\[
 v_x=\frac{b_0}{x}\,r_Rr_C\frac{u_D}{U_D}\frac{u_A}{U_A},
 \qquad xv_x\le C_\nu/\kappa_\nu^2 .
\]
\end{proposicion}
\begin{proof}
We have $U_f(1)=1$, $\kappa_\nu\le U_f(x)\le1$, and $dx/dt=-1$. Hence $\partial_t U_f=-m_f$.
Differentiating the product defining $J_x$ gives the second identity in \eqref{macro:derivadas}.
Substituting the factorization of $\nu_1$ gives the formula for $v_x$. The product identity and the
claims about the marginals follow by integration.
\end{proof}

The comparison that follows requires products of availability indicators for several lines.
Observing the two diagonal families separately gives only the means of the individual factors.
We now define sums of products and prove their concentration without assuming independence of the
lines.
